% from ega4-16-fr.tex
\oldpage[IV-4]{42}
est de même réunion de ses sous-groupes
$\operatorname{Diff}_{X/S}^n(\mathscr{F},\mathscr{G})$ (16.8.1). Les
structures de $\mathcal{O}_X$-Bimodule sur les
$\mathscr{D}\!iff_{X/S}^n(\mathscr{F},\mathscr{G})$ définissent donc une
structure de $\mathcal{O}_X$-Bimodule sur
$\mathscr{D}\!iff_{X/S}(\mathscr{F},\mathscr{G})$, explicitée encore par
(16.8.5.1) et (16.8.5.2).

Notons que, pour $n\leq m$, on a un diagramme commutatif
\[
  \xymatrix{
    \mathscr{H}\!om_{\mathcal{O}_X}
      (\mathscr{P}_{X/S}^n(\mathscr{F}),\mathscr{G})
      \ar[r]^-{\sim}\ar[d]
      &\mathscr{D}\!iff_{X/S}^n(\mathscr{F},\mathscr{G})\ar[d]\\
    \mathscr{H}\!om_{\mathcal{O}_X}
      (\mathscr{P}_{X/S}^m(\mathscr{F}),\mathscr{G})
      \ar[r]^-{\sim}
      &\mathscr{D}\!iff_{X/S}^m(\mathscr{F},\mathscr{G})
  }
  \tag{16.8.7.1}
  \label{IV.16.8.7.1}
\]
où les flèches horizontales sont les isomorphismes (16.8.4.2) et la flèche
verticale de gauche provient de l'homomorphisme canonique
$\mathscr{P}_{X/S}^m(\mathscr{F})\to
\mathscr{P}_{X/S}^n(\mathscr{F})$ (16.7.7). Pour tout ouvert $U$ de $X$,
munissons alors
\[
  \Gamma(U,\mathscr{P}_{X/S}^{\infty}(\mathscr{F}))
  =\varprojlim\Gamma(U,\mathscr{P}_{X/S}^n(\mathscr{F}))
\]
de la topologie limite projective des topologies discrètes sur les
$\Gamma(U,\mathscr{P}_{X/S}^n(\mathscr{F}))$, ce qui définit
$\Gamma(U,\mathscr{P}_{X/S}^{\infty}(\mathscr{F}))$ comme un
$\Gamma(U,\mathcal{O}_X)$-bimodule topologique, de sorte que
$\mathscr{P}_{X/S}^{\infty}(\mathscr{F})$ apparaît comme un faisceau à valeurs
dans la catégorie des groupes commutatifs topologiques
(0\textsubscript{I}, 3.2.6). Alors (G, II, 1.11) la limite du système inductif
de faisceaux de groupes commutatifs
\[
  (\mathscr{H}\!om_{\mathcal{O}_X}
  (\mathscr{P}_{X/S}^n(\mathscr{F}),\mathscr{G}))
\]
n'est autre que le faisceau des germes d'homomorphismes \emph{continus} de
$\mathscr{P}_{X/S}^{\infty}(\mathscr{F})$ dans $\mathscr{G}$ (ce dernier
étant muni de la topologie discrète): les homomorphismes continus de
$\Gamma(U,\mathscr{P}_{X/S}^{\infty}(\mathscr{F}))$ dans le groupe discret
$\Gamma(U,\mathscr{G})$ correspondent en effet biunivoquement aux systèmes
inductifs d'homomorphismes de groupes
$\Gamma(U,\mathscr{P}_{X/S}^n(\mathscr{F}))\to\Gamma(U,\mathscr{G})$. On peut
donc encore exprimer (16.8.4) en disant qu'on a un isomorphisme canonique
\[
  \mathscr{H}\!om.\mathrm{cont}_{\mathcal{O}_X}
  (\mathscr{P}_{X/S}^{\infty}(\mathscr{F}),\mathscr{G})
  \xrightarrow{\sim}
  \mathscr{D}\!iff_{X/S}(\mathscr{F},\mathscr{G})
\]
où le premier membre désigne le faisceau des germes d'homomorphismes continus
de $\mathscr{P}_{X/S}^{\infty}(\mathscr{F})$ dans $\mathscr{G}$.
\end{env}

\begin{proposition}[16.8.8]
\label{IV.16.8.8}
Soient $\mathscr{F}$, $\mathscr{G}$ deux $\mathcal{O}_X$-Modules,
$D:\mathscr{F}\to\mathscr{G}$ un homomorphisme de
$\psi^*(\mathcal{O}_S)$-Modules, $n$ un entier $\geq0$. Les conditions
suivantes sont équivalentes:
\begin{enumerate}
  \item[(a)] $D$ est un opérateur différentiel d'ordre $\leq n$.
  \item[(b)] Pour toute section $a$ de $\mathcal{O}_X$ au-dessus d'un ouvert
  $U$, l'homomorphisme $D_a:\mathscr{F}|U\to\mathscr{G}|U$ tel que, pour
  toute section $t$ de $\mathscr{F}$ au-dessus d'un ouvert $V\subset U$, on
  ait
  \[
    D_a(t)=D(at)-aD(t)
    \tag{16.8.8.1}
    \label{IV.16.8.8.1}
  \]
  est un opérateur différentiel d'ordre $\leq n-1$.
  \item[(c)] Pour tout ouvert $U$ de $X$, toute famille
  $(a_i)_{1\leq i\leq n+1}$ de $n+1$ sections de $\mathcal{O}_X$ au-dessus
  de $U$ et toute section $t$ de $\mathscr{F}$ au-dessus de $U$, on a
  l'identité
  \[
    \sum_{H\subset I_{n+1}}(-1)^{\operatorname{Card}(H)}
    \left(\prod_{i\in H}a_i\right)
    D\!\left(\left(\prod_{i\notin H}a_i\right)t\right)=0
    \tag{16.8.8.2}
    \label{IV.16.8.8.2}
  \]
  (où $I_{n+1}$ est l'intervalle $1\leq i\leq n+1$ de $\mathbb{N}$).
\end{enumerate}
